\section{Discussion}
\label{sec:discussion}

\textbf{Why the state comes apart.} The three tensors remember different things. The weights hold the whole run. The
second moment is a scale, averaged over $1/(1-\beta_2)$ steps, and it enters the update only through a square root, so
a relative error in it becomes a relative error half as large in the step. The first moment is a direction, averaged over
about ten steps. What it costs to leave a piece behind is set by two properties of that piece: how long it takes to
refill, and whether the step is too short or too long while it does. An empty first moment refills in ten steps and
shortens the step in the meantime. An empty second moment refills in a hundred and lengthens it. That asymmetry, and not
the byte count, orders Table~\ref{tab:main}: \arm{w\_v} and \arm{w\_m} send the same bytes and their peaks differ by a
factor of more than two hundred. It also explains the one strategy that pays for bytes it does not need. Carrying the
first moment next to a sketched second one, \arm{w\_m\_vr1}, lowers the peak from \ratioWVrone{} to \ratioWMVrone{} of
the band, and both are inside it, at twice the payload.

\textbf{Relation to earlier findings.} The conjecture of Section~\ref{sec:intro} holds at this scale: a run that could
have lived with a coarse second moment at every step~\cite{shazeer2018adafactor,dettmers2022eightbit} is not disturbed
by being handed one once. Two results qualify what is known. First, the error of the first restored update, the
acceptance test of Zhang et al.~\cite{zhang2026guarded}, is a safe screen but an incomplete one. \arm{w\_v} takes a
first step half the true size with a cosine of \cosV{} and would fail any tight bound on that error, yet it never
leaves the band, because the error is a shrinkage. \arm{w\_zero} and \arm{w\_fresh} take the same first step, \stepZero{}
and \stepFresh{} times too long, and end at \ratioWZero{} and \ratioWFresh{} bands, because what separates them happens
over the following steps. The direction of the error and its lifetime matter as much as its size. Second, a
weights-only restart is not one condition. Mitra~\cite{mitra2026orbittrace} reports that it costs final quality and that
rebuilding the first moment recovers little of the loss; here the five rules that send the weights alone span a factor
of nearly thirty in peak, and the rule that rebuilds the first moment is among the worst, for the reason given in
Section~\ref{sec:transient}. Our window ends \horizon{} steps after arrival, so we measure the transient and say
nothing about quality at the end of a long run.

\textbf{What a system can use.} Three things follow without any change to how a job trains. The changeover delay that
spot schedulers treat as a constant~\cite{wu2024cantbelate} can be replaced by a short list of pairs, bytes and
recovery cost, one per strategy, from which the scheduler picks inside the notice it has been given. The same payloads
serve as periodic checkpoints: the run of Section~\ref{sec:chaos} wrote its checkpoints with quantised moments, at half
the bytes of a full one. And the destination rule is free to change, since it costs no bytes. Loading weights into a
newly built optimizer is \arm{w\_fresh}; \warmPasses{} forward-backward passes before the first step turn it into
\arm{w\_warmv} and take the peak from \ratioWFresh{} bands to \ratioWWarmv. The method is reusable in the same way. A
bit-identical control and a seed band make an acceptance test for any lossy checkpoint, and both are cheap once data
order is a function of the step.

\textbf{Where we expect the picture to change.} These are predictions of the model and were not tested. At
$\beta_2 = 0.999$ the first update after a cold arrival is $\firstNoiseB$ times the steady one in place of
$\firstNoiseA$ and an empty second moment takes ten times as long to refill, so dropping both moments should cost more
and the sketch should be worth more. With a vocabulary of tens of thousands of tokens most embedding rows receive
gradient rarely, their second moments are far from a product of row and column sums, and the strategy that keeps them
exact may separate from the plain sketch, which it did not here. Early in training, at the peak learning rate, the
margin of the rank-one strategies is already thin (Section~\ref{sec:fork}). And the damage of dropping both moments
grew with model size across our range, from \sizeZeroXS{} to \sizeZeroM{} bands; we do not know whether it keeps
growing.
